\documentclass[book.tex]{subfiles}
\begin{document}

I was 11 years old when I wrote my first lines of software code. It happened on an MSX-1 computer in BASIC, and from that moment a completely new, magical world opened up to me. Five years later I got my first PC, an Intel 80286, and my fascination with computers only grew stronger. I witnessed the rise of the PC as a gaming platform and spent countless hours playing many of its groundbreaking titles: \textit{Prince of Persia}, \textit{Wolfenstein 3D}, \textit{Doom}, \textit{Dune II}, \textit{Command \& Conquer}, and many more.\\

\par
One of the first games I played on the PC was \textit{Commander Keen}. The moment I saw the game running, I was amazed. It was the first time I had seen smooth scrolling on a PC. Commander Keen was also the first major title from id Software. Thanks to the financial success of Commander Keen, the team---then still employed by Softdisk---decided to start their own game development company. The same studio would later go on to release the groundbreaking games \textit{Wolfenstein 3D} and \textit{Doom}.\\

\par
Fast forward to 2021, I discovered Fabien Sanglard’s website and began reading his Game Engine Black Books on Wolfenstein 3D and Doom. Inspired by those works, I wondered whether I could do something similar for Commander Keen: open up the source code, explore the files, and piece together a picture of the overall architecture and the clever tricks used. The style, dimensions, and structure of this book are intentionally similar to Fabien’s Game Engine Black Books, as an homage to those masterpieces. To give it a personal twist, I inverted the title and cover to white.\\

\par
The entire experience was immensely rewarding. It felt like returning to my early teenage years, debugging and experimenting with C and assembly code on an MS-DOS computer. I learned a great deal about the hardware of that era and the challenges developers faced in getting everything to work together.\\

\par
You might ask why I “wasted” my time on a game and hardware that are now more than 35 years old. Well, even though we now have vastly more transistors on a chip and have gone from megahertz to gigahertz CPUs, from kilobytes to gigabytes of RAM, and from simple video cards to powerful GPUs, the underlying architecture remains surprisingly similar to what it was three decades ago.\\


\par
The result of my journey is in this book---a mix of engineering, history, and nostalgia. It is written by a Dutch guy, so please forgive my imperfect English (ChatGPT was a great help in reviewing my grammar and style).  To compensate for my limitations in prose, I’ve included many detailed drawings to illustrate the hardware and software techniques discussed. All source code in this book is based on publicly available material. I have not altered the code, except for formatting and minor edits to improve clarity.\\

\par
So turn this page, sit back, and let’s travel together to the early ’90s.\\

\par
I hope you enjoy reading it.\\

-- Bas Smits\\
\par
Helmond, the Netherlands\\
\monthyeardate\today\\


\end{document}
